\section{二、理论分析与研究假说}\label{ux4e8cux7406ux8bbaux5206ux6790ux4e0eux7814ux7a76ux5047ux8bf4}

\subsection{（一）数据要素利用与资本市场定价效率}\label{ux4e00ux6570ux636eux8981ux7d20ux5229ux7528ux4e0eux8d44ux672cux5e02ux573aux5b9aux4ef7ux6548ux7387}

在剩余收益估值框架下，企业经济价值以账面净资产为锚（Ohlson，1995）。数据资产的大规模表外存在动摇了这一锚点：大量数据资源因不满足可辨认性或可靠计量条件而停留在报表之外，现行成本法又无法反映数据资源随应用场景扩展而非线性增值的特征（黄世忠等，2023），账面价值因此系统性低估了企业的真实经济价值。无形资产的定价偏差在近年日趋突出，中国A股市场的偏差尤为明显（Chen, J.等，2024；Dong和Doukas，2025；Bongaerts等，2025）。2024年施行的入表新规虽将数据资源纳入资产负债表，但适应期内投资者解读这类信息的难度并未立即下降。

这一计量缺口使市场参与者难以仅凭财务报表评估企业真实价值。企业在年度报告中披露的数据要素利用信息因此成为弥补计量缺口的增量信号，但这类信号对定价效率的影响方向并不确定。

信息供给效应强调信息不对称的改善作用。可信的公开信息披露压缩了知情与非知情交易者之间的信息差异，加速价格发现过程，这一机制既有理论上的严格论证（Grossman和Stiglitz，1980；Diamond和Verrecchia，1991），也有中国市场行业信息披露指引改善股价信息含量的独立实证支持（Gao等，2024）。信息模糊效应则指向相反方向，即信息不确定性越高，投资者越难以准确定价，价格发现因此越迟缓（Zhang，2006），数据驱动商业模式的内在复杂性、会计准则的估值空白以及数据丰裕条件下信息处理能力的制约（Dugast和Foucault，2025）都可能加剧这种不确定性。两种效应的行为假设不同，前者认为价格发现的约束在于信息供给不足，后者认为瓶颈在于投资者理解和处理信息的能力，净结果取决于哪种力量占主导。

在中国A股市场，若干制度条件可能有利于信息供给效应的发挥。年报是经过审计的法定文件，数据要素利用描述受到会计准则和审计程序的约束，其可验证性高于自愿性渠道。2022年``数据二十条''和2024年入表新规逐步降低了理解和评估数据资产信息的门槛，制度框架的渐进完善使信息模糊效应的制约有所减弱。若信息供给效应占优，应观察到企业数据要素利用与股价延迟呈显著负向关联。据此提出：

H1：企业数据要素利用降低了资本市场的股价延迟，即提升了定价效率。

\subsection{（二）定价效率的传导机制}\label{ux4e8cux5b9aux4ef7ux6548ux7387ux7684ux4f20ux5bfcux673aux5236}

上述假说关注的是总效应，但信息从企业年报向股票价格的具体传导路径尚待厘清。本文重点检验三条机制路径。

第一条是信息共识渠道。数据要素利用涉及技术架构、业务流程和价值实现方式，具有较强的专业门槛。如果年报披露能够提升外部投资者对企业数据利用方式的理解，分析师在处理相关信息时应当形成更一致的判断，预测分歧随之下降。分析师分歧的收敛意味着信息解释框架更趋一致，进而有助于相关信息更快融入股价。据此提出：

H2a：企业数据要素利用降低了分析师预测分歧。

第二条是经营稳定性渠道。数据驱动的运营管理有助于提升需求预测、库存配置和资源调度的效率，使企业经营活动更具可预测性。若数据要素利用缓解了经营层面的随机波动，企业现金流应更平稳，投资者面对的基本面不确定性也会下降，价格发现过程因而更顺畅。据此提出：

H2b：企业数据要素利用降低了现金流波动率。

第三条是供应链韧性渠道。供应链集中度过高意味着企业对少数关键客户或供应商依赖较强，一旦交易对手发生变化，基本面冲击更容易放大。数据要素利用可以帮助企业更准确地识别、筛选和匹配潜在交易伙伴，提高供应链管理效率并促进交易关系多元化。若这一过程成立，企业的综合供应链集中度应下降，基本面风险的暴露也将减弱。据此提出：

H2c：企业数据要素利用降低了综合供应链集中度。

\subsection{（三）效应的异质性}\label{ux4e09ux6548ux5e94ux7684ux5f02ux8d28ux6027}

总效应和传导机制之外，制度环境差异也可能调节定价效应的强弱。

第一，政策阶段差异影响数据利用信息的边际新颖性。相较于2020或2021，2022年``数据二十条''更接近数据要素制度密集推进的起点，也更适合作为主文中的阶段分界。在较早阶段，数据利用信息对投资者而言更具新颖性和辨识度，其边际定价价值可能更高；随着披露行为趋于制度化和常态化，边际信息增量可能减弱。据此提出：

H3a：企业数据要素利用对股价延迟的改善效应在2022年之前更强。

第二，信息中介活跃度可能影响数据利用披露的吸收速度。数据要素利用表述具有一定专业门槛，若分析师等信息中介对企业持续跟踪更充分，年报中的相关信息更容易被解释、传播并进入价格。相反，在分析师关注较低的企业中，即便存在披露，信息也可能因处理能力不足而更慢反映到价格中。据此提出：

H3b：企业数据要素利用对股价延迟的改善效应在高分析师关注企业中更强。

第三，上市板块差异可能改变数据利用披露的边际价值。主板公司通常经营模式更成熟、投资者覆盖更广、常规信息供给更充分，相关披露未必形成强增量信号；而非主板公司往往具有更强的成长性和更高的信息摩擦，数据利用披露可能更容易成为投资者重新理解企业价值的重要线索。据此提出：

H3c：企业数据要素利用对股价延迟的改善效应在非主板企业中更强。

产业属性差异同样可能影响市场的先验预期。例如，数字经济核心产业的企业天然与数据技术联系更紧密，相关披露未必构成强增量信息。但在统一滞后规格下，该类产业差异并未得到稳定支持，因此本文仅将其作为附录中的补充异质性结果报告，而不再据此建立主文的正式假说。
